% Part: first-order-logic
% Chapter: syntax-and-semantics
% Section: main-operator

\documentclass[../../../include/open-logic-section]{subfiles}

\begin{document}

\olfileid{fol}{syn}{mai}

\olsection{\printtoken{S}{main operator} یک فرمول}

\begin{explain}
اغلب سودمند است دربارهٔ آخرین عملگری سخن بگوییم که در ساخت
!!a{formula}~$!A$ به کار رفته است. این عملگر را \emph{عملگر اصلی}
فرمول~$!A$ می‌نامیم. به بیان شهودی، این عملگر «بیرونی‌ترین» عملگر $!A$ است.
برای نمونه، عملگر اصلی فرمول $\lnot !A$ نماد $\lnot$، و عملگر اصلی فرمول
$(!A \lor !B)$ نماد $\lor$ است، و به همین ترتیب.
\end{explain}


\begin{defn}[!!^{main operator}]
\ollabel{def:main-op}
\emph{!!{main operator}} !!a{formula}~$!A$ چنین تعریف می‌شود:
\begin{enumerate}
\item \indcase*{!A}{!A}{$\indfrm$ هیچ !!{main operator}ی ندارد.}

\tagitem{prvNot}{\indcase{!A}{\lnot !B}{!!{main operator} $\indfrm$
  نماد~$\lnot$ است.}}{}

\tagitem{prvAnd}{\indcase{!A}{(!B \land !C)}{!!{main operator}
  $\indfrm$ نماد~$\land$ است.}}{}

\tagitem{prvOr}{\indcase{!A}{(!B \lor !C)}{!!{main operator}
  $\indfrm$ نماد~$\lor$ است.}}{}

\tagitem{prvIf}{\indcase{!A}{(!B \lif !C)}{!!{main operator}
  $\indfrm$ نماد~$\lif$ است.}}{}

\tagitem{prvIff}{\indcase{!A}{(!B \liff !C)}{!!{main operator}
  $\indfrm$ نماد~$\liff$ است.}}{}

\tagitem{prvAll}{\indcase{!A}{\lforall[x][!B]}{!!{main operator}
  $\indfrm$ نماد~$\lforall$ است.}}{}

\tagitem{prvEx}{\indcase{!A}{\lexists[x][!B]}{!!{main operator}
  $\indfrm$ نماد~$\lexists$ است.}}{}
\end{enumerate}
\end{defn}

در هر حالت، منظورمان همان \emph{رخداد} مشخص و نشان‌داده‌شدهٔ
!!{main operator} در فرمول است. برای نمونه، چون فرمول
$((!D \lif !E) \lif (!E \lif !D))$ به صورت $(!B \lif !C)$ است که در آن
$!B$ برابر $(!D \lif !E)$ و $!C$ برابر $(!E \lif !D)$ است، رخداد دوم
$\lif$ همان !!{main operator} است.

\begin{explain}
این تعریف \emph{بازگشتیِ} تابعی است که همهٔ !!{formula}sی نااتمی را به
رخداد !!{main operator}شان می‌نگارد. به سبب شیوهٔ تعریف استقرایی
!!{formula}s، هر !!{formula}~$!A$ یکی از حالت‌های
\olref{def:main-op} را برآورده می‌کند. این امر تضمین می‌کند که برای هر
!!{formula} نااتمی~$!A$ یک !!a{main operator} وجود دارد. چون هر
!!{formula} فقط یکی از این شرط‌ها را برآورده می‌کند و چون در هر حالت
!!{formula}sی کوچک‌تری که $!A$ از آن‌ها ساخته شده به‌طور یکتا تعیین
می‌شوند، رخداد !!{main operator} فرمول~$!A$ یکتاست؛ پس تابعی تعریف کرده‌ایم.
\end{explain}

!!{formula}s را، بسته به اینکه !!{main operator} آن‌ها کدام نماد باشد، با
نام‌های \olref{tab:main-op} می‌خوانیم.\iftag{defNot,defOr,defAnd,defIf,defIff,defTrue,defFalse,defEx,defAll}
{ بااین‌حال، به یاد داشته باشید که عملگرهای تعریف‌شده رسماً در
!!{formula}s ظاهر نمی‌شوند. آن‌ها فقط اختصارند، پس رسماً نمی‌توانند عملگر
اصلی فرمولی باشند. ازاین‌رو در برهان‌های مربوط به همهٔ !!{formula}s لازم
نیست جداگانه بررسی شوند.}

\begin{table}[!h]
\centering
\begin{tabular}{c | c | c}
!!^{main operator} & نوع !!{formula} & نمونه\\
\hline
هیچ‌کدام & (!!{formula}) اتمی &
\iftag{prvFalse}{$\lfalse$،}{}
\iftag{prvTrue}{$\ltrue$،}{}
$\Atom{R}{t_1, \dots, t_n}$،
$\eq[t_1][t_2]$\\
$\lnot$ & نفی & $\lnot !A$ \\
$\land$ & عطف & $(!A \land !B)$ \\
$\lor$ & فصل & $(!A \lor !B)$ \\
$\lif$ & !!{conditional} & $(!A \lif !B)$ \\
$\liff$ & !!{biconditional} & $(!A \liff !B)$ \\
$\lforall[][]$ & (!!{formula}) کلی & $\lforall[x][!A]$ \\
$\lexists[][]$ & (!!{formula}) وجودی & $\lexists[x][!A]$
\end{tabular}
\caption{عملگر اصلی و نام‌های !!{formula}s}
\ollabel{tab:main-op}
\end{table}

\end{document}
